% !TeX program = lualatex
% halo:begin
% {
%   "version": 1,
%   "publish": true,
%   "course": "mathematical-analysis-i",
%   "section": "3.1",
%   "title": "数学分析（一）3.1：函数极限",
%   "categories": ["study-notes"],
%   "tags": ["mathematics"],
%   "excerpt": "函数在无穷远处、有限点处及单侧趋近时的极限定义，并用定义证明基本例题。"
% }
% halo:end
% 来源：IMG_2153.HEIC—IMG_2163.HEIC、IMG_2165.HEIC—IMG_2166.HEIC。
% 依照已有数列极限笔记的定义—证明—例题结构整理；端点处只从定义域内趋近。
\RequirePackage{iftex}
\RequireLuaTeX
\documentclass[UTF8,fontset=fandol,a4paper,12pt]{ctexart}
\usepackage[margin=2.5cm]{geometry}
\usepackage{amsmath,amssymb,amsthm}
\usepackage{enumitem,fancyhdr}
\usepackage[hidelinks]{hyperref}
\newcommand{\R}{\mathbb R}
\newcommand{\Nplus}{\mathbb N_{+}}
\newcommand{\U}{U}
\theoremstyle{definition}
\newtheorem*{definition}{定义}
\newtheorem{theorem}{定理}
\newtheorem*{criterion}{定理}
\renewcommand{\thetheorem}{3.\arabic{theorem}}
\newtheorem*{example}{例题}
\renewcommand{\proofname}{证明}
\setlist[enumerate]{leftmargin=2em,itemsep=0.2em,topsep=0.3em}
\setlength{\headheight}{16pt}
\pagestyle{fancy}
\fancyhf{}
\fancyhead[L]{\small 数学分析（一）}
\fancyhead[R]{\small 3.1\quad 函数极限}
\fancyfoot[C]{\thepage}
\renewcommand{\headrulewidth}{0.3pt}
\raggedbottom
\title{函数极限}
\author{Yuchen Zhang}
\date{整理日期：2026年10月14日}
\makeatletter
\renewcommand{\maketitle}{\begin{center}
{\LARGE\bfseries \@title\par}\vspace{0.6em}
\@author\par\vspace{0.2em}{\small\@date\par}
\end{center}}
\makeatother
\begin{document}
\maketitle
\setcounter{section}{2}
\section{函数极限}

数列极限研究下标趋于无穷时数列项的变化。函数极限研究的是：自变量按照指定方式趋近某个位置时，函数值是否趋近一个确定的数。这里的“趋近”只涉及函数定义域中可以取到的自变量，函数在趋近点的点值本身不决定极限。

\subsection*{自变量趋于无穷时的极限}
\begin{definition}
设函数 $f$ 在某个右半无限区间 $(a,+\infty)$ 上有定义。若存在 $A\in\R$，使得对任意 $\varepsilon>0$，都存在 $M>a$，当 $x>M$ 时有
\[
 |f(x)-A|<\varepsilon,
\]
则称 $f(x)$ 当 $x\to+\infty$ 时以 $A$ 为极限，记作
\[
 \lim_{x\to+\infty}f(x)=A.
\]
\end{definition}

类似地，若 $f$ 在某个左半无限区间 $(-\infty,a)$ 上有定义，且对任意 $\varepsilon>0$，存在 $M<a$，使 $x<M$ 时 $|f(x)-A|<\varepsilon$，则记作
\[
 \lim_{x\to-\infty}f(x)=A.
\]
若 $f$ 在 $|x|$ 足够大时有定义，并且
\[
 \lim_{x\to+\infty}f(x)=\lim_{x\to-\infty}f(x)=A,
\]
则称 $f(x)$ 当 $x\to\infty$ 时以 $A$ 为极限，记作
\[
 \lim_{x\to\infty}f(x)=A.
\]
此处的 $x\to\infty$ 表示 $|x|\to+\infty$，而不是只表示 $x\to+\infty$。几何上，若 $\lim_{x\to\infty}f(x)=A$，则直线 $y=A$ 是曲线的一条水平渐近线。

\begin{example}
证明 $\displaystyle\lim_{x\to+\infty}\sin\frac1x=0$。
\end{example}
\begin{proof}
任取 $\varepsilon>0$，取 $M=1/\varepsilon$。当 $x>M$ 时，利用
\[
 |\sin t|\le |t|\qquad(t\in\R),
\]
有
\[
 \left|\sin\frac1x-0\right|\le\frac1x<\frac1M=\varepsilon.
\]
故 $\displaystyle\lim_{x\to+\infty}\sin(1/x)=0$。
\end{proof}

\begin{example}
证明 $\displaystyle\lim_{x\to+\infty}\arctan x=\frac\pi2$。
\end{example}
\begin{proof}
任取 $\varepsilon>0$。当 $x>0$ 时，$\arctan x\in(0,\pi/2)$，故
\[
 0<\frac\pi2-\arctan x<\tan\left(\frac\pi2-\arctan x\right)=\frac1x.
\]
取 $M=1/\varepsilon$。当 $x>M$ 时，
\[
 \left|\arctan x-\frac\pi2\right|<\frac1x<\varepsilon.
\]
因此 $\displaystyle\lim_{x\to+\infty}\arctan x=\pi/2$。

由于 $\arctan x$ 是奇函数，同理可得
\[
 \lim_{x\to-\infty}\arctan x=-\frac\pi2.
\]
两个方向的极限不同，所以 $\arctan x$ 当 $x\to\infty$ 时不存在有限极限。
\end{proof}

\subsection*{自变量趋于有限点}
记
\[
 \U(x_0,\delta)=\{x\in\R:|x-x_0|<\delta\},\qquad
 \U^0(x_0,\delta)=\{x\in\R:0<|x-x_0|<\delta\}.
\]
\begin{definition}
设函数 $f$ 在 $x_0$ 的某个去心邻域 $\U^0(x_0,\delta_0)$ 上有定义。若存在 $A\in\R$，使得对任意 $\varepsilon>0$，存在 $\delta>0$，当
\[
 0<|x-x_0|<\delta
\]
时都有
\[
 |f(x)-A|<\varepsilon,
\]
则称 $f(x)$ 当 $x\to x_0$ 时以 $A$ 为极限，记作
\[
 \lim_{x\to x_0}f(x)=A.
\]
这里先给定 $\varepsilon$，再选取只依赖于 $\varepsilon$ 和 $x_0$ 的 $\delta$；选定后，同一个去心邻域内的所有 $x$ 都必须满足误差估计。由于条件排除了 $x=x_0$，所以 $f(x_0)$ 可以没有定义，也可以任意指定，而不影响该极限。
其中可取 $0<\delta\le\delta_0$。
\end{definition}

若极限定义中的 $\delta$ 可以取为 $\delta_0$，那么任意 $0<\delta\le\delta_0$ 也可以作为一个选择。证明中常把 $\delta$ 取得更小，以同时满足几个条件。

\subsection*{单侧极限}
若只允许 $x$ 从 $x_0$ 的左侧趋近，记
\[
 \lim_{x\to x_0^-}f(x)=A;
\]
若只允许 $x$ 从右侧趋近，记
\[
 \lim_{x\to x_0^+}f(x)=A.
\]
例如，右极限的定义是：对任意 $\varepsilon>0$，存在 $\delta>0$，使得当 $x_0<x<x_0+\delta$ 时，$|f(x)-A|<\varepsilon$；左极限同理，将条件改为 $x_0-\delta<x<x_0$。

\begin{criterion}[双侧极限的判定]
\[
 \lim_{x\to x_0}f(x)=A
 \quad\Longleftrightarrow\quad
 \lim_{x\to x_0^-}f(x)=\lim_{x\to x_0^+}f(x)=A.
\]
\end{criterion}
\begin{proof}
必要性显然：双侧去心邻域包含左侧和右侧的去心邻域，因此双侧极限成立时，两侧极限都等于 $A$。

反之，设左右极限都等于 $A$。任取 $\varepsilon>0$，分别取 $\delta_1,\delta_2>0$，使左侧和右侧的误差均小于 $\varepsilon$。令 $\delta=\min\{\delta_1,\delta_2\}$。当 $0<|x-x_0|<\delta$ 时，$x$ 必在 $x_0$ 左侧或右侧，因而总有 $|f(x)-A|<\varepsilon$。故双侧极限等于 $A$。
\end{proof}

\begin{example}
设
\[
 f(x)=\begin{cases}x+2,&x\ge0,\\x-2,&x<0.\end{cases}
\]
则
\[
 \lim_{x\to0^+}f(x)=2,\qquad \lim_{x\to0^-}f(x)=-2.
\]
左右极限不相等，故 $\lim_{x\to0}f(x)$ 不存在。
\end{example}

\begin{example}
函数 $f(x)=|x|/x$ 在 $x=0$ 的去心邻域内有定义。因为
\[
 \lim_{x\to0^-}\frac{|x|}{x}=-1,qquad
 \lim_{x\to0^+}\frac{|x|}{x}=1,
\]
所以 $\lim_{x\to0}|x|/x$ 不存在。
\end{example}

\newpage
\subsection*{用定义证明极限}
证明有限点处的极限时，先把 $|f(x)-A|$ 化成可控制的表达式；若表达式含有随 $x$ 变化的分母，先用 $|x-x_0|<1$ 等条件使分母远离零。

\begin{example}
证明
\[
 \lim_{x\to2}\frac{x^2-4}{x-2}=4,
 \qquad
 \lim_{x\to2}\frac{x-2}{x^2-4}=\frac14.
\]
\end{example}
\begin{proof}
当 $x\ne2$ 时，$x^2-4=(x-2)(x+2)$。

（1）任取 $\varepsilon>0$，取 $\delta=\varepsilon$。当 $0<|x-2|<\delta$ 时，
\[
 \left|\frac{x^2-4}{x-2}-4\right|=|x+2-4|=|x-2|<\varepsilon.
\]

（2）先要求 $|x-2|<1$，于是 $1<x<3$，从而 $|x+2|>3$。任取 $\varepsilon>0$，取
\[
 \delta=\min\left\{1,\,12\varepsilon\right\}.
\]
当 $0<|x-2|<\delta$ 时，
\[
 \left|\frac{x-2}{x^2-4}-\frac14\right|
 =\left|\frac1{x+2}-\frac14\right|
 =\frac{|x-2|}{4|x+2|}
 <\frac{|x-2|}{12}<\varepsilon.
\]
故两式均得证。
\end{proof}

\begin{example}
证明对任意 $x_0\in\R$，有
\[
 \lim_{x\to x_0}\sin x=\sin x_0.
\]
\end{example}
\begin{proof}
由和差化积公式，
\[
 |\sin x-\sin x_0|
 =2\left|\sin\frac{x-x_0}{2}\right|\left|\cos\frac{x+x_0}{2}\right|
 \le 2\left|\sin\frac{x-x_0}{2}\right|
 \le |x-x_0|.
\]
任取 $\varepsilon>0$，取 $\delta=\varepsilon$。当 $0<|x-x_0|<\delta$ 时，上式给出
\[
 |\sin x-\sin x_0|<\varepsilon.
\]
\end{proof}

同理，利用 $\cos x-\cos x_0=-2\sin((x+x_0)/2)\sin((x-x_0)/2)$，可得
\[
 \lim_{x\to x_0}\cos x=\cos x_0.
\]

\begin{example}
证明
\[
 \lim_{x\to1^-}\sqrt{1-x^2}=0,
 \qquad
 \lim_{x\to-1^+}\sqrt{1-x^2}=0.
\]
\end{example}
\begin{proof}
函数 $\sqrt{1-x^2}$ 的定义域是 $[-1,1]$，所以在端点处只能从定义域内的一侧趋近。

以 $x\to1^-$ 为例。任取 $\varepsilon>0$，取 $\delta=\min\{1,\varepsilon^2/2\}$。当 $1-\delta<x<1$ 时，
\[
 0<1-x^2=(1-x)(1+x)<2(1-x)<2\delta\le\varepsilon^2,
\]
因而 $|\sqrt{1-x^2}-0|<\varepsilon$。另一端点的证明完全类似，故结论成立。
\end{proof}

\end{document}
